\section{What Survives: Coherence-vs-Similarity Disentanglement}
\label{sec:coherence_noise}

The previous three sections weakened the original case. This
section reports the one structural finding that does survive
scrutiny: at matched per-passage similarity, coherence-preserving
uninformative noise produces a smaller faithfulness drop than random
noise.

\paragraph{Why this matters.}
Section \ref{sec:causal_intervention} showed that the specific
operationalization of coherence as
$\mathrm{CCS} = \overline{\mathrm{sim}} - \sigma_{\mathrm{sim}}$
does not predict faithfulness at fixed mean similarity. That null
is interesting only if a different, milder claim --- that
coherence is a real variable separable from similarity --- is also
testable.

\paragraph{Protocol.}
For each SQuAD query (\(n{=}200\), seed 42), we retrieve top-3
passages and run four conditions:
\begin{itemize}
\itemsep -2pt
\item \textbf{Baseline}: top-3 retrieval, no perturbation.
\item \textbf{Random off-topic noise}: replace 1, 2, or 3 of the 3
retrievals with random off-topic passages from the same corpus.
\item \textbf{Coherence-preserving uninformative noise}:
replace 1, 2, or 3 of the 3 retrievals with passages from the
\emph{same query's top-20 retrieval pool} that do not contain the
answer span. Same-topic, just not answer-bearing.
\item \textbf{HCPC-v1 refinement}: the published refinement
operator that motivated this paper.
\end{itemize}
For each condition we measure mean retrieval similarity to the
query and mean DeBERTa-v3 NLI faithfulness. The slope of
faithfulness as a function of noise rate is the regression slope of
faith on noise rate $\in\{0, 1/3, 2/3, 1\}$.

\paragraph{Result.}

\begin{table}[H]
\centering
\small
\caption{Coherence-preserving noise vs.\ random noise on SQuAD,
$n{=}200$. Slopes are linear regression of faithfulness on noise rate.
``Drop at full noise'' is faith at noise=$1$ minus faith at
noise=$0$.}
\label{tab:coherence-noise}
\begin{tabular}{lccc}
\toprule
Condition & Faith.\ slope per noise rate &
Sim.\ slope per noise rate & Drop at full noise \\
\midrule
Random off-topic noise         & $-0.069$ & $-0.481$ & $0.052$ \\
\textbf{Coherence-preserving} (same-topic, answer-absent) &
$\mathbf{-0.043}$ & $\mathbf{-0.113}$ & $0.042$ \\
HCPC-v1 refinement (post-hoc) & --- & --- & $0.001$ \\
\bottomrule
\end{tabular}
\end{table}

The slopes are markedly different. Random noise drives faithfulness
down by $0.069$ per unit noise rate while also crashing similarity
($-0.481$). The coherence-preserving condition --- in which the
\emph{same} number of passages is replaced, but with same-topic
passages that don't carry the answer --- drops faithfulness by only
$0.043$ per unit noise, with similarity barely moving ($-0.113$).

\paragraph{Interpretation.}
At matched per-passage similarity, structural coherence carries a
real-but-modest faithfulness signal. Random noise removes both
similarity and coherence simultaneously; the coherence-preserving
condition isolates coherence and finds a smaller, but non-zero,
faithfulness drop. The ratio of slopes (random $-0.069$ over
coherence-preserving $-0.043$) is $\sim 1.6\times$, well within
the noise of any single experiment but consistent across noise
rates.

Combined with the Section \ref{sec:causal_intervention} null, the
honest statement is: \emph{coherence operationalized as the CCS
score does not predict faithfulness at fixed mean similarity, but
some structural property of coherent retrieval --- not yet
correctly captured by any score we tested --- carries signal
beyond similarity.}

\paragraph{This is not a re-rescue.}
We do not claim that the original refinement-paradox claim is
restored by this finding. The SQuAD/Mistral-7B paradox magnitude
collapsed at scale (\S\ref{sec:scaling}) and its proposed
mechanism (\S\ref{sec:causal_intervention}) was disconfirmed.
What \S\ref{sec:coherence_noise} adds is a directional and
empirically robust observation that ``random noise'' is not the
right null for retrieval-quality experiments. Future work
operationalizing coherence in a way that survives the
matched-similarity test of \S\ref{sec:causal_intervention} would
inherit this section's structural support.
